\begin{abstract}
\end{abstract}

\section{Case Studies}
\label{casestudies}

\subsection{Feature-based Multivariate Time Series Analysis}
Feature-based time series analysis \cite{}
% feature-based methods are invariant to the dimension of the temporal component of a time series, and so enable statistical and machine learning techniques to draw comparisons between time series based on their relationships in a common feature space. in multivariate time series, the spatial dimension ($M$ for an M-channeled multivariate process) may also vary, and so feature-based methods -- while eliminating temporal scale variance -- still result in a feature set which differs in dimension due to the spatial component, and so cannot be put into a common space, and therefore cannot be analyzed with conventional machine learning models. "Chopping" datasets to the lowest common dimension will often eliminate important dynamical regimes \text{it} between channels of a multivariate system, and thereby erode important information about the underlying system interactions. What are the current methods to circumvent this problem of domain shift variance (I.e., many processes in a multivariate time series) while still preserving the important dynamical properties and dependency structures in the entire system?

CONTEXT: For an $M$-channeled multivariate time series (MTS) $\mathbf{X}^{M \times T}$ with channels $\{X_t^{(i)}\}$ for $i \in \{1, 2, \dots, M\}$ the python library \texttt{pyspi} computes a diverse library of $>250$ statistics for pairwise interaction (SPIs) between every pair of channels, yielding ${M \choose 2}=\frac{M(M-1)}{2}$ total interactions for a chosen statistic. Computing the $k$-th SPI yields an $(M \times M)$-dimensional adjacency matrix, coined a matrix of pairwise interaction (MPI). Within this adjacency matrix, we can geometrically interpret the $i$-th node to represent the $i$-th channel of the multivariate time series, $X_t^{(i)}$, where the elements $(i,j)$ are the edge-weights corresponding to the value of the SPI computed between the $i$-th and $j$-th nodes, i.e., $\mathrm{SPI}_k(i,j)$. By computing $K$ SPIs on a given multivariate time series, we generate a $(K \times M \times M)$-dimensional feature space, with each $\mathrm{MPI}_k$ the adjacency matrix corresponding to the $k$-th SPI.

Like feature-based univariate time series analysis, operating in MPI-space allows us to compare multivariate time series which may differ in the length of the time component. For example, let us jointly consider two MTS processes: $\mathbf{X}^{M \times T_1}$ and $\mathbf{Y}^{M \times T_2}$, where $T_1 \neq T_2$. Then, computing the $k$-th SPI for each pair of channels within $\mathbf{X}$ and $\mathbf{Y}$ yields the two adjacency matrices $\mathrm{MPI}_k({\mathbf{X}})$ and $\mathrm{MPI}_k({\mathbf{Y}})$. Comparing the matrix element $(i,j)$ between each MPI informs on the strength of the chosen statistic between nodes (channels) $i$ and $j$ within each time series. This reveals information on the dependency structure between channels within each system according to the choice of statistic. For example, if we choose to compute the Pearson correlation ($r$) (the SPI) between every pair of channels within each system, we can read the values of the resultant MPI to identify which nodes are emblematic of a linear dependency structure (high $(i,j)$), or those which are indicative of a non-Gaussian relationship, such as nonlinear variation or heavy-tailed outliers.

In a similar way that feature-based methods for univariate time series are helpful to compare (potentially varying length) data using a set of summary statistics (features), a strength of using \texttt{pyspi} to convert a multivariate time series into a $K$-length series of interpretable features (MPIs) is its invariance to datasets differing in their temporal scale. This is useful to convert multivariate time series data into a feature space sharing a common dimension -- the number of selected features $k$ -- which can, in turn, be used for comparing systems varying in sampling rate or length with more explanatory statistical or machine learning techniques. However, like all feature-based methods for multivariate systems, they remain sensitive to the size of the spatial component of the process -- the number of channels $M$ -- which limits the breadth of analysis between systems. Thus, if we could compare systems which vary in potentially both the spatial and temporal component, we should be able to uncover relationships between multivariate systems which may be dynamically, dependently, or structurally related in non-trivial ways.

\subsection{Extracting Informative Features from Multivariate Time Series}
A possible novel route to circumvent this is to construct an abstracted feature space generated from the relationships \textit{between} features themselves. In a similar style to which feature-based methods for univariate time series abstract one dimension from the data (they transform a $T$-dimensional data into a uniform, interpretable $K$-dimensional feature space), abstracting a second dimension resolves the dependency on the spatial component, $M$. 

% For an undirected statistic -- mutual information, cross-correlation, coherence magnitude -- this results in $M(M-1)/2$ unique elements; for a directed statistic -- transfer entropy, time-lagged mutual information, phase slope index -- we have $M(M-1)$ unique elements. 
Consider a set of $K$ SPIs and two MTS processes, $\mathbf{X}^{M_1 \times T_1}$ and $\mathbf{Y}^{M_2 \times T_2}$, where $M_1 \neq M_2$ (and optionally $T_1 \neq T_2$). Computing the $K$ SPIs on each MTS yields $K$ MPI adjacency matrices, each with $M(M-1)$ off-diagonal elements $(i, j)$. By extracting the off-diagonal entries of each MPI and computing the correlation \textcolor{red}{(I think Pearson over Spearman (evaluate the choices) because rank-space may corrupt subtle signatures between clusters of pairs of interactions between channels)} between the entries of these matrices yields a singular summary statistics -- a feature $f_k \coloneqq \mathrm{corr}(\mathrm{SPI}_{\mathbf{X}}, \mathrm{SPI}_{\mathbf{Y}})$. For $K$ SPIs, the result is a ${K \choose 2}$-dimensional feature vector, where each value encodes the empirical relationship between any two statistics for a given multivariate time series. Since this feature vector is uniform in its dimensionality for any two multivariate systems $\mathbf{X}$ and $\mathbf{Y}$ -- irrespective of the counts spatial node count $M$ or the temporal length $T$ -- it uniquely allows for \textit{any} set of multivariate systems to be embedded into a common feature space, which is necessary to perform downstream statistical and machine learning.

For consistency, we call this feature space generated ``\textit{SPI-SPI space}". Crucially, in this abstracted feature space, the information extracted from each value is interpreted differently to conventional feature-based analyses. Comparing the values of this feature vector between different multivariate processes informs not on the strength of relationship between channels, but on the empirical similarity in the \textit{character} or \textit{nature} of dependency structures within each system. For example, consider two multivariate systems, of any spatial and temporal construction, and a series of $K$ SPIs whose intersection of statistical assumptions can highlight the presence (or absence) of a particular dependency regime present in each system. In the first system, define a strong coupling of some characteristic nature between nodes $(1, 2)$, and the same coupling between nodes $(1, 3)$ in the second system. Since the ``\textit{meshgrid}" of statistical assumptions is assumed to be capable of detecting the dynamical regime underpinning the coupling, the relationships between pairs of the $K$ statistics (the features $f_k$ of the SPI-SPI feature space $\mathbf{f}^{K \choose 2}$) are sufficient to express the presence (or absence) of the coupling regime within the system. That is, when considered jointly, the relevant features $f_k \in \mathbf{f}$ will co-vary in such a way that we can conclude that both systems exhibit the toy coupling as defined in each of the two systems. We note that since this feature space is a further dimension abstract from the raw data, it is agnostic to some \textit{fundamental} differences within the system. That is, while this method is capable of detecting the shared dependency structure within systems (given a sufficiently expressible set of statistics), it is ``blind`` to \textit{where} the coupling lives. Hence, this feature space is most suitable to analyses which involve \textcolor{red}{[unfinished].}


\subsection{Case Studies}
We can establish the intuition of this method through an examination of the intricate dynamical signatures that are yielded in carefully constructed multivariate systems using this feature-feature method.

\subsubsection{Differentiating Between Linear, Nonlinear, and Non-monotonic Dependencies}
In this study, we investigate how the feature-feature correlation space can tease a linear or nonlinear dependency within multivariate systems.

Consider Pearson’s product-moment correlation coefficient $r$ and Spearman's rank-correlation coefficient $\rho$, which are two of the most well-known statistics for quantifying the strength and direction of pairwise dependence between time series. In many settings, the results of their computation yield similar results, since Spearman's $\rho$ is equivalent to Pearson’s $r$ computed on rank-transformed data. Thus, for datasets where variables are distributed linearly in the absence of non-Gaussian outliers, we should expect reasonable equivalence (i.e., a strongly coupled $f_k = \mathrm{corr}(r, \rho)$). Contrastingly, for systems that include heavy-tailed data or strongly non-Gaussian scaling -- like seismograph readings in earthquakes or financial returns with heavy-tailed outliers -- Pearson’s $r$ and Spearman’s $\rho$ will begin to \textit{meaningfully} differ, since Spearman’s rank-transformed correlation space is invariant to large monotonic scaling changes or data polluted with non-Gaussian outliers. In the case of the decoupling of $r$ with $\rho$ (i.e., a relative decrease in $f_k \coloneqq \mathrm{corr}(r, \rho)$ compared to our linear system above), we could conclude that our data must have some characteristic nonlinearity to it that is not captured by the assumptions of Pearson correlation. That is, our data must exhibit at least one of (monotonic) non-linear scaling, heavy-tailed outliers, or non-Gaussian distributed variables. [unfinished]